\subsection{Subject-Namespace Claims}
\label{sec:namespace-claims}

The results so far hold for any admissibility function satisfying
Assumption~\ref{as:disjoint}. Two further questions need more structure: when a
sub-query whose patterns have variable subjects can be delegated, and how an engine
settles claim status in practice. Both rely on claims that are decided by the subject of
a triple alone.

\begin{definition}[Subject-determined claims]
\label{def:subject-determined}
Claims are \emph{subject-determined} if for every composite $c$ there is a set of terms
$S_c$ such that $\alpha_c(t) = 1$ iff $\mathrm{subj}(t) \in S_c$.
\end{definition}

For this work we instantiate admissibility by namespaces, as discovery-based querying on
the Web identifies sources by IRI. With $\mathrm{ns}(c)$ the namespace of the IRI of $c$,
\[
  \rho_c(t) = 1 \iff \mathrm{subj}(t) \in \mathrm{ns}(c),
\]
and a composite that declares claims its whole namespace, $\sigma_c = \rho_c$. Claims are
then subject-determined with $S_c = \mathrm{ns}(c)$. A composite with namespace
\url{http://example.org/} thus claims every triple whose subject starts with it.
$\rho_c$ depends only on the identity of $c$, so the engine knows it as soon as it
discovers $c$, whatever lies behind the frontier. Admissibility bounds what $c$ may
\emph{claim}, not what it may \emph{aggregate}; for the latter, write
$\mathrm{iri}(d)$ for the IRI of a source $d$.

\begin{assumption}[Namespace aggregation]
\label{as:ns-aggregation}
A composite aggregates every source in its namespace: if
$\mathrm{iri}(d) \in \mathrm{ns}(c)$ then $d \in \Dd_c$.~\todo{align with subweb paper}
\end{assumption}

The picture is that of horizontal partitioning, with each composite the sole source of
record for a region of the subject space. The difference is that the partitioning is
declared by the participants rather than imposed by a coordinator, so no global catalogue
of it exists and both the participants and the partitioning are discovered during
execution. Assumption~\ref{as:disjoint} holds when no composite's namespace extends
another's. Disjoint aggregation alone does not suffice: a composite on
\url{http://example.org/} and one on \url{http://example.org/a/} may aggregate disjoint
documents and still both claim a triple about \url{http://example.org/a/x}.

\subsubsection{Resolving claim status}
\label{sec:claim-status}

Call $t \in \dataf(d)$ \emph{local} if no composite namespace separates its subject from
its source: for every composite $c$, $\mathrm{subj}(t) \in \mathrm{ns}(c)$ implies
$\mathrm{iri}(d) \in \mathrm{ns}(c)$. Under Assumption~\ref{as:ns-aggregation} a local
triple is covered, since a composite claiming it has $\mathrm{subj}(t) \in \mathrm{ns}(c)$,
hence $d \in \Dd_c$ and $t \in \dataf(c)$. By Lemma~\ref{lem:covered-stable}, solutions
built from local triples are sound as soon as they are produced; only solutions using a
non-local triple, a statement about a resource in another namespace, can be invalidated
by later discovery.

\begin{remark}[Resolving claim status by request]
\label{rem:claim-lookup}
In practice the engine need not wait for discovery to settle whether a non-local triple
$t$ is claimed. By Assumption~\ref{as:disjoint} at most one composite can claim $t$, and
under namespace admissibility its location is fixed by $\mathrm{subj}(t)$ alone: it is the
composite, if any, whose namespace contains $\mathrm{subj}(t)$. The engine can therefore
resolve the claim status of $t$ with a request to this possible aggregator, for example
by looking up an advertised composite along the path of $\mathrm{subj}(t)$ and, if one
claims $t$, asking it whether it holds $t$. Including these lookups in $\disc$ makes the
aggregator part of $\reach_B(W)$, so the outcome agrees with the answer under claims. An
engine that resolves the status of every non-local triple before using it emits only
solutions that remain valid, at the cost of one lookup per foreign namespace and one
membership check per non-local triple, rather than a wait for the frontier to empty.
This mirrors authoritative reasoning, where authority is decided per term by
dereferencing~\cite{hogan2009authoritative}.~\todo{Mechanism for advertising a
composite, and for establishing that none exists: storage description, link header?}
\end{remark}

\subsubsection{Star- and chain-shaped sub-queries}
\label{sec:shapes}

Subject-determined claims make the claim status of a matched triple a function of the
binding of its subject. For a star whose patterns share the subject variable $?s$, all
triples matched by a solution $\mu$ have subject $\mu(?s)$, so either $c$ claims all of
them or none. The star is therefore claim-uniform for every composite, whatever the data,
and Theorem~\ref{th:claim-delegation} applies at any time $T$; $\Phi_c$ keeps the
solutions with $\mu(?s) \in S_c$. Patterns with distinct subject variables, as in chains,
admit mixed solutions.~\todo{Formalise for stars and size-2
chains}

\todo{Section indicating how anchor variables are required to correctly filter out partial results}
\todo{This then makes stars and size 2 chains attractive candidates to reduce client work, while
longer chains can't really be pruned.}
